% Cuerpo científico recuperado y distribuido en su residencia terminal.
\section{Cofinal Closure of the Riemann--Weil Criterion}

The genealogical realization produces a cofinal, natural family of
positive forms whose publication coincides with the complete
Guinand--Weil functional. Extinction of the terminal and compatibility with the
inclusions permit passage to the limit without introducing a list of zeros or the
sought sign. By Weil's reversible criterion, every nontrivial zero
\(\rho\) of the Riemann zeta function satisfies
\[
  \operatorname{Re}\rho=\frac12.
\]
The regulated finite-window computations are subsequent controls; they do not form
part of the premises of the cofinal closure.

\paragraph{Cofinal Factorization of the Spectral Conclusion.}
The probative chain gathers, in the already established causal order, the finite
genealogical form, its positive complement, naturality under inclusions, and
compatible extinction of the terminal term. Each window preserves the prime
clocks, the Archimedean channel, the orientation of the two sheets, and the boundary
register. Cofinality jointly transports those data to the limit, and the
Guinand--Weil identity identifies the resulting form with its classical
spectral publication.

Consequently, positivity does not proceed from a subsequent selection of
zeros: it descends from the positive operator constructed prior to evaluation. The
functional involution preserves the mirror sheet and fixes the line
\(\operatorname{Re}\rho=1/2\); cofinal compatibility preserves that fixation
in the passage from regulated windows to the complete form. The finite
calculations thus perform their proper function of subsequent verification, whereas
mathematical closure resides in the preceding positive and natural composition.

